%=============================================================================!
% 16.7  CHARGE TRANSPORT                                                      !
%=============================================================================!
\section{Charge transport}
\label{sec:ob-charge}

An electrolyte carries current through the motion of its ions. If each
ion moved independently, its diffusion coefficient would determine its
contribution to the conductivity. But a cation and an anion that move
together as a neutral pair can diffuse without carrying net charge.
We first derive the independent-ion estimate, then include correlations
and measure their effect in a model molten salt.

\subsection*{Drift against diffusion}

Suppose each ion of some kind feels a steady force $F_x$ along $x$, for
instance from an electric field. Its collisions with its neighbours act,
on average, as a friction that grows with its velocity, so it does not
speed up for ever: it drifts with a mean velocity $v_{\mathrm d}$ in the
direction of the force, at which the friction balances the force, while
diffusion spreads it. The drift and the diffusion are tied together. If
the force comes from a potential energy $U(x)$, so that $F_x = -\dd U/\dd
x$, the ions settle into the canonical density $\rho(x) \propto
\mathrm{e}^{-\beta U(x)}$ (\cref{sec:sm-canonical}), whose slope is
$\dd\rho/\dd x = -\beta\rho\,\dd U/\dd x = \beta F_x\rho$. At equilibrium no
net flow remains, so the drift's flow $\rho\,v_{\mathrm d}$ and Fick's flow
$-D\,\dd\rho/\dd x$ cancel:
\begin{equation*}
   \rho\,v_{\mathrm d} - D\beta F_x\rho = 0 , \qquad
   v_{\mathrm d} = \beta D\,F_x ,
\end{equation*}
the Einstein relation: the drift velocity per unit force is the
diffusion coefficient divided by $\kB T$. The same atom that wanders
fast also drifts fast when pushed.

In an electric field $\mathcal{E}$ along $x$, ion $i$, of charge $q_i$,
feels $F_x = q_i\mathcal{E}$ and drifts at $\beta D_iq_i\mathcal{E}$, with
$D_i$ the diffusion coefficient of its kind. The current density, the
charge carried across unit area in unit time, is the sum over the ions of
each charge times its drift velocity, divided by the volume, $\sum_iq_i
\times \beta D_iq_i\mathcal{E}/V$. Divided by $\mathcal{E}$, this is the
conductivity of independent ions, the Nernst-Einstein value
\begin{equation}
   \sigma_{\mathrm{NE}} = \frac{1}{V\kB T}\sum_iq_i^2D_i ;
   \label{eq:ob-ne}
\end{equation}
cations and anions add, $q_i^2$ being positive for both. In the units of
\texttt{mdlab}, $e^2$ per \si{\angstrom}, \si{\femto\second} and
\si{\electronvolt}, a conductivity of 1 is \SI{1.602e6}{\siemens\per\metre},
the siemens being the ampere per volt.

\subsection*{The conductivity of the whole}

Ions do not move independently. The general result has the form of
\cref{eq:ob-einstein} for the charge displacement $\vb{M}_q(t) =
\sum_iq_i\vb{r}_i(t)$, the sum of every ion's position weighted by its
charge:
\begin{equation}
   \sigma_{\mathrm e} = \frac{1}{6V\kB T}\lim_{t\to\infty}\frac{\dd}{\dd t}
   \ens{\lvert\vb{M}_q(t_0 + t) - \vb{M}_q(t_0)\rvert^2} ,
   \label{eq:ob-sigma}
\end{equation}
whose derivation, like the viscosity's, uses the response of the system
to a weak field\cite{AllenTildesley2017}. Expanding the square shows how
it differs from \cref{eq:ob-ne}. With $\Delta\vb{r}_i$ the displacement of
ion $i$,
\begin{equation*}
   \ens{\lvert\Delta\vb{M}_q\rvert^2}
   = \sum_iq_i^2\ens{\lvert\Delta\vb{r}_i\rvert^2}
   + \sum_i\sum_{j\neq i}q_iq_j\ens{\Delta\vb{r}_i\cdot\Delta\vb{r}_j} .
\end{equation*}
The first sum grows as $\sum_iq_i^2\,6D_it$ and gives
\cref{eq:ob-ne}; the second holds the correlations between different
ions. A cation and an anion that move together have $q_iq_j < 0$ and
$\ens{\Delta\vb{r}_i\cdot\Delta\vb{r}_j} > 0$, a negative term, and lower
the conductivity below $\sigma_{\mathrm{NE}}$. Their ratio is the Haven
ratio $H_{\mathrm R} = \sigma_{\mathrm{NE}}/\sigma_{\mathrm e}$. Like $D$,
$\sigma_{\mathrm e}$ also follows from a Green-Kubo integral, of the
autocorrelation function of the charge current $\sum_iq_i\vb{v}_i$, by
the steps of \cref{eq:ob-msd-vacf}.

\subsection*{A model molten salt}

The test system is 64 ions in a cubic cell of side \SI{13.4}{\angstrom},
32 of charge $+e$ with the mass of sodium and 32 of $-e$ with that of
chlorine. Every pair repels as $A\mathrm{e}^{-r/r_{\mathrm{rep}}}$ with
$r_{\mathrm{rep}} = \SI{0.3}{\angstrom}$, and $A = \SI{6231}{\electronvolt}$ is
set so that a pair of opposite charges has its least energy at
\SI{2.8}{\angstrom}, where the repulsion's slope $A\mathrm{e}^{-r/r_{\mathrm
{rep}}}/r_{\mathrm{rep}}$ balances the
Coulomb pull $k_{\mathrm e}e^2/r^2$; the Coulomb energy is summed by Ewald's
method (\cref{sec:md-ewald}). The model is illustrative, chosen to
resemble an alkali halide, a salt of a metal such as sodium with a
halogen such as chlorine, and is not fitted to sodium chloride. Eight
runs each melt the lattice at \SI{4000}{\kelvin}, hold
\SI{1500}{\kelvin} with CSVR for \SI{40}{\pico\second} and then run
\SI{1}{\nano\second} at fixed energy with steps of \SI{2}{\femto\second}.
Their energies differ, since 64 ions leave the thermostatted stage with an
energy that fluctuates widely, so they sit at temperatures from 1362 to
\SI{1721}{\kelvin}.

The liquid is ordered by charge (\cref{fig:ob-salt}a). Around a cation
the anions crowd at \SI{3.07}{\angstrom}, where $g_{+-}$ rises to 3.65,
and other cations lie farther out, $g_{++}$ peaking at 1.66 at
\SI{4.37}{\angstrom}, $g_{--}$ alike at \SI{4.32}{\angstrom}; each cation
has 6.76 anions out to the first minimum of $g_{+-}$, at
\SI{4.87}{\angstrom}.

Each run gives $\sigma_{\mathrm e}$ from \cref{eq:ob-sigma}, with the
charge displacement recorded at every step and fitted over 1 to
\SI{5}{\pico\second}, and $\sigma_{\mathrm{NE}}$ from the two kinds' $D$
over the same window. The Green-Kubo route agrees with the displacement
in every run to within 0.6\%. Both conductivities rise with the run's
temperature (\cref{fig:ob-salt}b), and the Nernst-Einstein value lies
above the true one by more than half a percent in six runs of eight,
within it in one and below it in one. Over the eight, at $(1487 \pm
45)$\,\si{\kelvin}, $\sigma_{\mathrm e} = (158 \pm 9)$\,\si{\siemens\per
\metre} and $\sigma_{\mathrm{NE}} = (168 \pm 10)$\,\si{\siemens\per\metre},
errors that come mostly from the spread of the runs' temperatures. Their
ratio, taken run by run, shows no trend with the temperature in the
figure, and averages
\begin{equation*}
   H_{\mathrm R} = 1.062 \pm 0.022 ,
\end{equation*}
2.8 errors above 1, beyond the 2.365 that Student's factor for eight
runs, seven degrees of freedom, sets at 95\% (\cref{sec:cv-mean}): the
ions' motions are correlated, and treating them as independent
overstates the conductivity of this melt by about 6\%. The diffusion
coefficients of the two kinds over the eight runs, $(5.19 \pm 0.46)$ and
$(5.03 \pm 0.45)\times10^{-4}$\,\si{\angstrom\squared\per\femto\second}
for the cations and the anions, are the same within errors that again
reflect the temperatures; taken run by run, which removes the
temperature, their ratio is $1.033 \pm 0.008$, the cations faster in
seven runs of eight.

\begin{figure}[htb]
\centering
\includegraphics{ch16_observables/salt}
\caption{The model molten salt, 64 ions in eight runs of
\SI{1}{\nano\second} at fixed energy. (a) The partial radial distribution
functions $g_{+-}$, $g_{++}$ and $g_{--}$. (b) The conductivity of each
run from the charge displacement (circles) and from Nernst-Einstein with
the ions' own diffusion coefficients (squares), against the run's
temperature. (c) The vibrational density of states of the cations and of
the anions in the first run (\cref{sec:ob-vdos}).}
\label{fig:ob-salt}
\end{figure}

\begin{keyidea}
An ion pushed by a force $F_x$ drifts at $\beta DF_x$, so independent ions
conduct as $\sigma_{\mathrm{NE}} = \sum_iq_i^2D_i/V\kB T$. The true
conductivity comes from the mean squared displacement of the total charge
displacement $\sum_iq_i\vb{r}_i$, whose cross terms hold the correlations
between ions; the Haven ratio $\sigma_{\mathrm{NE}}/\sigma_{\mathrm e}$
measures them.
\end{keyidea}

\notebook{16}{compare the charge displacement with the ions' own displacements}

\begin{selfcheck}
If every cation were permanently bound to one anion, moving as a neutral
pair, what would $\sigma_{\mathrm e}$ be, and what would
$\sigma_{\mathrm{NE}}$ still give?
\end{selfcheck}
